\chapter{Conclusions and Future Work}\label{ch:conclusion}

\section{Summary of Contributions}\label{sec:concl_summary}
This thesis developed and evaluated a modular high-performance
software stack for scalable and efficient Large Language Models. The
principal contributions are:
\begin{enumerate}
    \item A comparative study of four distributed training strategies
          on an 8-GPU server, with detailed utilisation, memory, and
          energy measurements (Chapter~\ref{ch:distributed}).
          DeepSpeed ZeRO-3 attained the highest peak GPU utilisation
          ($93.8\%$) and the best per-epoch energy ($22\%$ over DP);
          FSDP matched ZeRO-3 on memory efficiency while remaining
          simpler to deploy.
    \item A treatment of CUTLASS/CuTe GEMM pipelines and their
          mapping onto transformer layers
          (Chapter~\ref{ch:cutlass}). CUTLASS matched cuBLAS on large
          square GEMMs and outperformed it by $35\%$ on skinny
          attention shapes and $31\%$ on grouped MoE shapes; epilogue
          fusion saved an additional $10$--$25\%$ of end-to-end HBM
          traffic.
    \item A novel Sliding Window Flash Attention algorithm with block
          classification, vectorised mask generation, up to $65.6\%$
          block skipping, and a closed-form optimal-window formula
          whose loss is strictly convex
          (Chapter~\ref{ch:swfa}). SWFA enabled $50$--$59\%$
          throughput improvement on GPT-2-Medium training and
          achieved sub-$10^{-4}$ FP32 accuracy over $19$ validation
          configurations.
    \item Kernel-level optimisations inside the SGLang serving
          runtime for long-context inference
          (Chapter~\ref{ch:sglang}), delivering a cumulative
          $2.84\times$ decode-throughput speedup on LLaMA-7B with
          $w = 4096$ while keeping p99 latency under $30$\,ms at
          batch $64$.
    \item A Triton-based profiling and auto-tuning pipeline
          (Chapter~\ref{ch:triton}) that reproduced the CUTLASS
          numbers to within $3\%$, providing an independent
          cross-check on the SWFA and SGLang measurements.
\end{enumerate}
The end-to-end evaluation shows a $50$--$59\%$ throughput improvement
and $25$--$29\%$ energy reduction for a GPT-2-Medium workload when
SWFA is combined with DeepSpeed or FSDP, demonstrating that the
individual layers of the stack compose rather than compete.

\section{Lessons Learned}\label{sec:concl_lessons}

\subsection{Layering Wins Over Point Optimisation}\label{sec:concl_layering}
No single optimisation in this thesis delivers the end-to-end result
in isolation: the $59\%$ training throughput gain requires the SWFA
kernel, the sharded training schedule, and the epilogue-fused GEMMs
to cooperate. The SGLang chapter makes the same point from a
different angle---the $2.84\times$ serving speedup is the product of
four independent kernels, each responsible for a distinct bottleneck.

\subsection{Block-Level Sparsity Is Practical}\label{sec:concl_sparsity}
Sparsity at the element level is hard to exploit on modern Tensor
Cores: the $16 \times 8 \times 16$ MMA instructions issue whole dense
tiles. Block-level sparsity, in contrast, integrates cleanly because
each tile is either fully present or fully absent. SWFA illustrates
this at the scale of the attention matrix, and the same pattern should
extend to other structured sparsity families (strided, blocksparse,
BigBird).

\subsection{Profiling Is a First-Class
Deliverable}\label{sec:concl_profiling}
Having a uniform profiling pipeline (Chapter~\ref{ch:triton}) turned
out to be as valuable as the kernels themselves: it allowed every
number in the thesis to be reproduced with a single command and
flagged several false positives (most notably the $20\%$ regression
that turned out to be a miskeyed auto-tune cache entry).

\section{Limitations}\label{sec:concl_limitations}

\subsection{Single-Node Scope}\label{sec:concl_single_node}
All measurements are on a single 8-GPU A100 node. Multi-node scaling
would introduce an InfiniBand (or similar) hop whose latency
characteristics are fundamentally different from NVLink-3. Strategies
that minimise cross-node traffic (pipeline parallelism, 2-D sharding,
gradient compression) would likely shift the recommendation away from
pure ZeRO-3.

\subsection{A100 Focus}\label{sec:concl_a100}
The kernel numbers are specific to A100. H100 brings asynchronous
WGMMA, TMA, DSMEM, and an FP8 transformer engine, all of which shift
the knee of the roofline and would require retuning. The CUTLASS and
Triton code paths are structurally ready for H100; only the
auto-tune search space and the tile shapes need updating.

\subsection{Prototype-Level SWFA
Kernel}\label{sec:concl_swfa_proto}
The SWFA kernel is implemented in PyTorch with CPU-optimised paths.
A production CUDA/Triton kernel is sketched in
Chapter~\ref{ch:triton} but has not yet been deployed in an
end-to-end serving run. The $50$--$59\%$ training speedup numbers are
therefore likely a lower bound; a hand-tuned CUTLASS port should push
further.

\section{Future Work}\label{sec:concl_future}

\subsection{Short term (6--12 months)}\label{sec:concl_short}
\begin{itemize}
    \item Complete the CUDA/Triton port of SWFA and its backward
          pass; target a further $2$--$3\times$ speedup via Tensor-Core
          utilisation. Extend the block classifier to handle
          user-declared sparse patterns (strided, global, BigBird-like)
          within the same kernel skeleton.
    \item Finish the Triton auto-tuner campaign on H100 and publish a
          shape-indexed configuration cache that covers the most
          common LLM shapes out of the box.
    \item Integrate SWFA into the SGLang prefill kernel so that the
          decoupling between the prefill and decode attention paths
          is complete.
\end{itemize}

\subsection{Medium term (1--2 years)}\label{sec:concl_medium}
\begin{itemize}
    \item Multi-vendor support: port the stack to Intel Xe GPUs via
          oneDNN and SYCL, and to AMD via ROCm and Composable Kernel.
          The CuTe layout algebra is the portable core; the vendor-%
          specific parts are the MMA instruction traits.
    \item Hierarchical communication for multi-node training, with
          gradient compression and error compensation. A first step is
          to integrate Zero++ with InfiniBand-aware NCCL algorithms.
    \item Automatic policy selection: a lightweight runtime that
          chooses between DDP/FSDP/ZeRO-3 and the best attention
          kernel based on the model shape and the available hardware.
          The scaling-laws literature~\cite{hoffmann2022chinchilla}
          suggests that the decision can be precomputed offline from
          $(M, N, d, w, P)$ alone.
\end{itemize}

\subsection{Longer term}\label{sec:concl_long}
\begin{itemize}
    \item Co-design of the attention pattern with the training curriculum:
          the cost model of Chapter~\ref{ch:swfa} assumes a fixed
          attention decay $\lambda$, but in practice $\lambda$ changes
          during training. A schedule that grows $w$ as $\lambda$
          decreases would turn $w^\star$ into a time-varying curriculum.
    \item Integration with emerging low-precision formats (FP4, NVFP6)
          that require fused scale/dequantisation epilogues; both
          CUTLASS and Triton already support the necessary fused
          instructions on Hopper.
    \item Carbon-aware scheduling of training runs: combining the
          energy-per-epoch numbers of Chapter~\ref{ch:distributed}
          with a real-time grid carbon signal would shift training
          windows to low-carbon hours without extending the wall-clock
          of the campaign.
\end{itemize}

\section{Closing Remarks}\label{sec:concl_closing}
Efficient training and serving of LLMs is not the product of a single
algorithmic trick---it requires simultaneously tuning the distributed
schedule, the kernel implementation, and the attention algorithm. The
work presented here approaches each of these layers in isolation and
then shows that they compose, yielding substantial improvements in
throughput, memory use, and energy efficiency while remaining
numerically faithful. As models continue to grow, this holistic
view---rather than point optimisations within a single framework---is
likely to remain the right way to deliver practical performance on
the hardware that is actually available today.

The principal hope of the author is that the pieces delivered in this
thesis, most importantly the SWFA algorithm with its proof of
correctness, the controlled distributed-training comparison, and the
Triton profiling pipeline, serve as building blocks that are useful
well beyond the specific workloads evaluated here. The full algorithmic
pseudocode and reference PyTorch implementations are given in
Appendix~A; the rigorous proofs underlying the main theoretical
results are in Appendix~B.
